\documentclass[10pt,a4paper]{amsart}
\usepackage[margin=19mm]{geometry}
\usepackage{amsmath,amssymb,mathtools}
\usepackage[scr=rsfs,cal=euler]{mathalfa}
\usepackage{xfrac}
\usepackage[unicode,hidelinks,bookmarks=false]{hyperref}
\newcommand{\ie}{i.e.\ }
\newcommand{\cf}{cf.\ }
\newcommand{\resp}{respectively\ }
\newcommand{\Subsection}{\subsection}
\providecommand{\nobreakdash}{\nobreak\hbox{-}}
\setlength{\emergencystretch}{2em}
\multlinegap0pt
\usepackage{bm}
\usepackage{amssymb}
\usepackage{mathtools}
\usepackage{enumerate}
\usepackage{xcolor}
\usepackage{algorithm}\usepackage{algpseudocode}
\newtheorem{theorem}{Theorem}
\newcommand{\Norm}[1]{\|#1\|}
\newcommand{\R}{\mathbb{R}}
\def\eqref#1{Eq-(\ref{#1})}
\def\gO{{\mathcal{O}}}
\def\gX{{\mathcal{X}}}
\def\mA{{\bm{A}}}
\def\mJ{{\bm{J}}}
\def\mU{{\bm{U}}}
\def\mV{{\bm{V}}}
\def\mW{{\bm{W}}}
\def\vu{{\bm{u}}}
\def\vw{{\bm{w}}}
\def\vx{{\bm{x}}}
\title{Towards A Unified PAC-Bayesian Framework for Norm-based Generalization Bounds}
\author{Xinping Yi, Gaojie Jin, Xiaowei Huang,, Shi Jin}
\date{Source: arXiv:2601.08100v1 (CC BY 4.0)}
\begin{document}
\maketitle

\noindent\emph{Source and licence.} Towards A Unified PAC-Bayesian Framework for Norm-based Generalization Bounds. Version arXiv:2601.08100v1 (CC BY 4.0), distributed by arXiv under the Creative Commons Attribution 4.0 International licence, http://creativecommons.org/licenses/by/4.0/, as recorded in the arXiv licence metadata for this exact version and shown on the abstract page https://arxiv.org/abs/2601.08100v1. Full licence notice: \texttt{licenses/arxiv-2601.08100-CC-BY-4.0.txt}.\par\smallskip

\noindent\textit{Editorial scope.} Counted unit: the article's theorem thm:low-rank-matrix (source0.tex lines 1046--1060). The article's complete original proof of it is delivered with it (source0.tex lines 1061--1063, one \texttt{proof} environment of 3 lines). No mathematics is written, repaired or re-derived: the statement and the proof are the article's own lines, in the article's own notation, and no retained line is substituted or re-flowed.  Retained uncounted as the source-local context the proof needs: lines 705--707; lines 1319--1340.  Nothing was omitted from any retained span, so the package carries no omission record.  No retained line contains a citation, so the wrapper carries no bibliography and no bibliography entry.  No retained line contains a figure environment, an image-inclusion command or drawing code, so no image is delivered and the package declares no asset dependency.  Editorial formatting only, in addition: an AMS \texttt{amsart} preamble that restates the article's own declarations and macros used by the retained lines, namely the environments \texttt{theorem} on separate counters, together with the standard packages those lines need.  The retained environments print in this document as Theorem 1 in order of appearance, whereas the article prints the same items with the numbering of its own full document.  The complete native source of the article as distributed in the arXiv e-print for this exact version is bundled unchanged as \texttt{sources/192575/source0.tex}.
\begin{align} \label{eq:pert-bound-origin}
    \Norm{f_{\vw+\vu}(\vx)-f_{\vw}(\vx)}_2 \le e B \prod_{l=1}^d \|\mW_l\|_2 \sum_{l=1}^d \frac{\|\mU_l\|_2}{\|\mW_l\|_2}
\end{align}

\begin{align}
    L_0(f_{\vw}) \le \hat{L}_{\gamma}(f_{\vw}) + \mathcal{O} \left( \sqrt{ \frac{B^2d^2h^2 \Phi^{\mathrm{rn}}(\vw) + \ln \frac{md}{\delta}}{\gamma^2m}} \right)
\end{align}
for any $\beta$. 
This completes the proof.



\subsection{Proof of Theorem \ref{thm:low-rank-matrix}}
\label{proof:low-rank-matrix}
The procedure is similar to that in the diagonal case, while the difference lies in the following two points:
\begin{itemize}
    \item Without relying on the perturbation bound in \eqref{eq:pert-bound-origin}, the difference of the network's output due to weight perturbation is approximated by its first-order Taylor expansion with the Jacobian $\mJ_l(\vx)$.
    \item The sensitivity matrix $\mA_l$ captures the directions aligning with the right singular vectors of $\mJ_l(\vx)$, so that the scaling factor of generalization bound reduces from $h^2$ to $K$.
\end{itemize}
\subsubsection{Perturbation condition}
The perturbation difference is the change in the output $f_{\vw}(\vx)$ when the weights are perturbed by $\vu$. It is influenced by the Jacobian matrix $\mJ_l(\vx)$ at each layer, which determines how sensitive the output is to perturbations in the weights.

By the first-order Taylor expansion of $f_{\vw}(\vx)$ with a small perturbation $\vu$, we have
\begin{align}
    f_{\vw+\vu}(\vx) - f_{\vw}(\vx) = \sum_{l=1}^d \mJ_l(\vx) \vu_l  + o(\Norm{\vu_l}_2)
\end{align}

\begin{theorem}\label{thm:low-rank-matrix}
Consider a feedforward neural network $f_{\vw}: \gX \to \R^K$ with ReLU activations, depth $d$, width $h$, bounded inputs $\Norm{\vx}_2 \le B$, and a potential low-rank structure.
    By setting a low-rank sensitivity matrix as 
\begin{align} \label{eq:low-rank-Al}
    \mA_l = \sqrt{d} \mV_l \mathrm{diag}(\sigma_{l,1},\dots,\sigma_{l,K},0,\dots,0)\mV_l^T
\end{align}
with $K\ll h^2$, where $\mV_l$ is the collection of right eigenvectors of the Jacobian of the output with respect to the weights at the $l$-th layer, i.e., $\mJ_l(\vx) = \frac{\partial f_{\vw}(\vx)}{\partial \mathrm{vec}(\mW_l)} \in \R^{K \times h^2}$, and
\begin{align}
    \sigma_{l,k} = B\frac{\prod_{l=1}^d \Norm{\mW_l}_2}{\Norm{\mW_l}_2}, \ \forall~k,
\end{align}
 for any $\delta,\gamma>0$, with probability at least $1-\delta$, for any $\vw$, the generalization error over a training set with size $m$ can be upper-bounded by
    \begin{align}
    L_0(f_{\vw}) \le \hat{L}_{\gamma}(f_{\vw}) + \gO \left( \sqrt{ \frac{B^2d^2K \Phi(\vw) + \ln \frac{dm}{\delta}}{\gamma^2m}} \right).
    \end{align}
\end{theorem}

\begin{proof}
See Appendix \ref{proof:low-rank-matrix}.
\end{proof}

\end{document}
